%%%%%%%%%%%%%%%%%%%%%%%%%%%%%%%%%%%%%%%%%%%%%%%%%%%%%%%%%%%%%%%%
\section{结论与讨论}
%%%%%%%%%%%%%%%%%%%%%%%%%%%%%%%%%%%%%%%%%%%%%%%%%%%%%%%%%%%%%%%%

\subsection{主要发现}

本研究并未发现一个跨模型普遍成立的“LLM代理风险偏移”。
更稳妥的结论是：LLM风险偏好具有显著模型异质性，并会被代理角色、问责线索、损益框架、推理模式和题目呈现方式共同塑造。

\textbf{H1A（算法谨慎偏移）}获得有限支持。
收益域内，DeepSeek-V4-Flash和DeepSeek-V4-Pro方向上更谨慎，但BH校正后不显著；GPT-4.1仅在可观察代理条件下显著谨慎化。
跨损益域合并的补充检验显示，DeepSeek-V4-Pro存在轻度总体谨慎偏移（$q_{BH}=0.035$）。
因此，H1A应被理解为模型特定、条件性的发现。

\textbf{H1B（道德风险偏移）}未获总体支持，但在GPT-4.1收益域不可观察代理条件下获得局部支持（$q_{BH}=0.019$）。
这说明低问责代理语境可能触发风险放松，但该模式没有在其他模型上稳定复现。

\textbf{H2（问责性与损益域调节）}仅在GPT-4.1上获得条件性支持。
差中差检验显示，GPT-4.1中可见性带来的谨慎化在损失域强于收益域（DiD=+1.22，$q_{BH}=0.025$）；其他模型未呈现显著交互。

\textbf{H3（模型异质性）}获得强支持。
Grok-4在整卷作答中接近风险中性，Kimi-k2高度保守，GPT-4.1与DeepSeek V4系列更具情境敏感性。
附加检验还表明，深度思考模式和逐题独立调用会改变部分模型的绝对风险水平。

\subsection{理论与实践意义}

理论上，本文将替他人决策和前景理论扩展到LLM代理行为评估中。
结果显示，LLM并不简单复制人类风险偏好：多数模型在损失域反而更趋向风险规避，提示安全对齐、风险语义和显式推理过程可能共同重塑模型选择。
这对前景理论的AI外推具有提醒意义。
人类的损失域风险寻求通常被解释为损失厌恶、翻盘动机和主观价值函数弯曲的共同结果；而LLM没有真实财富账户和情绪性损失体验，它对“损失”的反应更可能来自训练语料中风险规避话语、对齐规则和任务提示。
因此，本文并不是否定前景理论，而是说明在合成被试研究中，理论概念需要重新映射到模型可观察行为上。

实践上，AI金融代理不能被默认视为风险中性工具。
金融应用上线前应进行模型内、任务内的行为审计，并关注透明审计、推理模式和题目呈现方式等调控变量。
尤其在投顾、保险和信贷等场景中，开发者不能只测试模型是否会给出“合理解释”，还需要测试其在不同责任归属、可见性和损益框架下是否系统性改变建议。
如果某一模型在低问责条件下更愿意承担风险，平台需要通过日志审计、用户确认和风险提示降低委托人受损的概率；如果模型在损失语义下过度保守，则可能需要在产品层面区分风险控制和机会成本。

\subsection{对实验设计核心问题的回答}

本研究问题是：LLM在担任金融代理人、替他人进行风险决策时，其风险偏好是否相对于为自己决策发生偏移，以及这种偏移是否受到损益域和问责压力调节。
这一问题具有理论意义，因为它连接了人类代理决策、前景理论和LLM行为偏差；也具有现实意义，因为LLM正在进入投顾和交易辅助等高风险场景。

实验方法适合回答该问题。
Holt-Laury量表可以将风险偏好转化为安全选择数，2 $\times$ 3设计能够区分代理角色、可观察性和损益框架。
Self、Agent-Invisible和Agent-Visible分别提供基线、低问责代理情境和可观察问责情境，也符合课程对AI synthetic respondents的要求。

\subsection{研究局限与未来方向}

本研究仍有六点限制。

（1）\textbf{提示词语义污染。}
实验材料中将选项A标注为“安全选项”、选项B标注为“风险选项”，这可能直接诱导LLM响应“安全/风险”语义，而不仅仅基于概率和金额做选择。
此外，不可观察代理条件中“委托人不关心结果”的表述同时改变了问责性和委托人偏好强度。
因此，本文结论应被理解为在当前提示词语境下成立；未来需要加入中性标签和“委托人关心但不可观察”的稳健性实验。

（2）\textbf{API重复调用不等同于人类独立被试。}
每个条件下的50次调用是同一模型的独立生成样本，而不是50个具有独立偏好的人类个体。
统计检验反映的是模型在给定Prompt下的输出分布差异。

（3）\textbf{实验情境的外部有效性有限。}
Holt-Laury彩票范式有较高内部效度，但与真实金融代理中的资产配置、组合管理和多轮反馈仍有距离~\cite{Eriksen2010Myopic}。

（4）\textbf{模型和参数覆盖有限。}
本文仅测试5款模型，且使用各平台默认采样参数。
未来应系统操纵temperature、top-p、推理开关和模型版本。

（5）\textbf{题目呈现方式本身是实验因素。}
逐题独立调用会改变部分模型的绝对风险水平，未来应将整卷/逐题作为正式因子，而不仅作为稳健性检查。

（6）\textbf{缺少人类被试对照。}
本文主要与既有文献比较，尚未在同一材料、同一统计口径下采集人类样本。
引入人类对照组将有助于直接判断人机代理决策的异同。
